\subsection{GRADED RINGS AND MODULES}

DEFINITION 2. Given a ring A and a graduation \((\mathbf{A}_{\lambda})\) of type \(\Delta\) on the additive group A, this graduation is said to be compatible with the ring structure on A if

\[
\mathrm{A} _ {\lambda} \mathrm{A} _ {\mu} \subset \mathrm{A} _ {\lambda + \mu} \quad f o r a l l \lambda , \mu i n \Delta .\tag{1}
\]

The ring A with this graduation is then called a graded ring of type \(\Delta\) .

PROPOSITION 1. If every element of \(\Delta\) is cancellable and \((\mathbf{A}_{\lambda})\) is a graduation of type \(\Delta\) compatible with the structure of a ring \(\mathbf{A}\) , \(\mathbf{A}_0\) is a subring of \(\mathbf{A}\) (and in particular \(1 \in \mathbf{A}_0\) ).

As \(\mathbf{A}_0\mathbf{A}_0\subset \mathbf{A}_0\) by definition, it suffices to prove that \(1\in \mathbf{A}_0\) . Let \(1 = \sum_{\lambda \in \Delta}e_{\lambda}\) be the decomposition of 1 into its homogeneous components. If \(x\in \mathbf{A}_{\mu}\) , then \(x = x.1 = \sum_{\lambda \in \Delta}xe_{\lambda}\) ; comparing the components of degree \(\mu\) , (since \(\mu +\lambda = \mu\) implies \(\lambda = 0\) ) \(x = xe_0\) . Since this relation is true for every homogeneous element of A, it is true for all \(x\in \mathbf{A}\) ; in particular \(1 = 1.e_0 = e_0\in \mathbf{A}_0\) .

DEFINITION 3. Let A be a graded ring of type \(\Delta\) , \((\mathbf{A}_{\lambda})\) its graduation and M a left (resp. right) A-module; a graduation \((\mathbf{M}_{\lambda})\) of type \(\lambda\) on the additive group M is compatible with the A-module structure on M if

\[
\mathbf {A} _ {\lambda} \mathbf {M} _ {\mu} \subset \mathbf {M} _ {\lambda + \mu} \quad (\text { resp. } \mathbf {M} _ {\mu} \mathbf {A} _ {\lambda} \subset \mathbf {M} _ {\lambda + \mu})\tag{2}
\]

for all \(\lambda, \mu\) in \(\Delta\) . The module \(\mathbf{M}\) with this graduation is then called a left (resp. right) graded module of type \(\Delta\) over the graded ring \(\mathbf{A}\) .

When the elements of \(\Delta\) are cancellable, it follows from (2) and Proposition 1 that the \(\mathbf{M}_{\lambda}\) are \(\mathbf{A}_0\) -modules.

Clearly if A is a graded ring of type \(\Delta\) , the left A-module \(A_{s}\) (resp. the right A-module \(A_{d}\) ) is graded of type \(\Delta\) .

Examples. (1) On any ring A the trivial graduation of type \(\Delta\) is compatible with the ring structure. If A is graded by the trivial graduation, for a graduation ( \(M_{\lambda}\) ) of type \(\Delta\) on an A-module M to be compatible with the A-module structure, it is necessary and sufficient that the \(M_{\lambda}\) be submodules of M.

\begin{enumerate}
\def\labelenumi{(\arabic{enumi})}
\setcounter{enumi}{1}
\tightlist
\item
  Let A be a graded ring of type \(\Delta\) , M a graded A-module of type \(\Delta\) and \(\rho\) a homomorphism of \(\Delta\) into a commutative monoid \(\Delta'\) whose identity element is denoted by 0. Then A is a graded ring of type \(\Delta'\) and M a graded module of type \(\Delta'\) for the graduations of type \(\Delta'\) derived from \(\rho\) and the graduations of type \(\Delta\) on A and M by the procedure of no. 1, Example 1: this follows immediately from the relation \(\rho(\lambda + \mu) = \rho(\lambda) + \rho(\mu)\) .
\end{enumerate}

In particular, if \(\Delta = \Delta_{1} \times \Delta_{2}\) is a product of two commutative monoids, the projections \(pr_{1}\) and \(pr_{2}\) are homomorphisms and the corresponding graduations are just the partial graduations derived from the graduations of type \(\Delta\) (no. 1, Example 3); these partial graduations are thus compatible with the ring structure of A and the module structure of M.

Similarly, if \(\Delta = \Delta_0^{(I)}\) (where \(\Delta_0\) is a commutative monoid with identity element denoted by 0), the total graduation (no. 1, Example 4) of type \(\Delta_0\) derived from the graduation of type \(\Delta\) on A (resp. M) by means of the codiagonal homomorphism is compatible with the ring structure on A (resp. with the module structure on M).

\begin{enumerate}
\def\labelenumi{(\arabic{enumi})}
\setcounter{enumi}{2}
\tightlist
\item
  Let A be a graded ring of type \(\Delta\) , M a graded A-module of type \(\Delta\) and \(\lambda_0\) an element of \(\Delta\) ; for all \(\lambda \in \Delta\) , let \(M_{\lambda}' = M_{\lambda + \lambda_0}\) and let \(M'\) be the Z-module \(\bigoplus_{\lambda \in \Delta} M_{\lambda}'\) . As \(A_{\lambda}M_{\mu}' \subset M_{\lambda + \mu + \lambda_0} = M_{\lambda + \mu}'\) , \(M'\) is an A-module and the \(M_{\lambda}'\) form on M' a graduation of type \(\Delta\) compatible with the A-module structure of M'; the graded A-module M' of type \(\Delta\) thus defined is said to be obtained by shifting by \(\lambda_{0}\) the graduation of M and it is denoted by \(\mathbf{M}(\lambda_{0})\) . When \(\Delta\) is a group, the underlying A-module of the graded A-module M' is identified with M.
\end{enumerate}

*(4) Let B be a commutative ring. The polynomial ring B{[}X{]} in one indeterminate is graded of type N by the subgroups BX \(^{n}\) (n ≥ 0) (cf.~III, § 2, no. 9 and IV).*

*(5) Let B be a commutative ring, E a B-module, Q a quadratic form on E and C(Q) the Clifford algebra of Q (cf.~IX, §9). The sub-B-modules C \(^{+}\) (Q) and C \(^{-}\) (Q) form on C(Q) a graduation of type Z/2Z compatible with the ring structure on C(Q).*

Remarks. (1) The graduations most often used are of type Z or of type \(Z^{n}\) ; when we speak of graded (resp. bigraded, trigraded, etc.) modules and rings without mentioning the type, it is understood that we mean graduations of type Z (resp. \(Z^{2}\) , \(Z^{3}\) , etc.); a graded ring (resp. module) of type N is also called a graded ring (resp. module) with positive degrees.

\begin{enumerate}
\def\labelenumi{(\arabic{enumi})}
\setcounter{enumi}{1}
\tightlist
\item
  The graded \(\mathbf{Z}\) -modules of type \(\Delta\) , when \(\mathbf{Z}\) has the trivial graduation, are just the graded commutative groups (whose set of degrees is a commutative monoid) of Definition 1 (no. 1).
\end{enumerate}

DEFINITION 4. Let \(\mathbf{A}, \mathbf{A}'\) be two graded rings of the same type \(\Delta\) and \((\mathbf{A}_{\lambda})\) , \((\mathbf{A}_{\lambda}')\) their respective graduations. A ring homomorphism \(h: \mathbf{A} \to \mathbf{A}'\) is called graded if \(h(\mathbf{A}_{\lambda}) \subset \mathbf{A}_{\lambda}'\) for all \(\lambda \in \Delta\) .

Let \(\mathbf{M}, \mathbf{M}'\) be two graded modules of type \(\Delta\) over a graded ring of type \(\Delta\) . Let \(u: \mathbf{M} \to \mathbf{M}'\) be an A-homomorphism and \(\delta\) an element of \(\Delta\) ; \(u\) is called graded of degree \(\delta\) if \(u(\mathbf{M}_{\lambda}) \subset \mathbf{M}_{\lambda + \delta}\) for all \(\lambda \in \Delta\) .

Let A be a graded ring of type \(\Delta\) , \(A'\) a graded ring of type \(\Delta'\) and \(\rho: \Delta \to \Delta'\) a homomorphism. A ring homomorphism \(h: A \to A'\) is called graded if \(h\) is a graded homomorphism of graded rings of type \(\Delta'\) when A is given the graduation of type \(\Delta'\) derived from its graduation of type \(\Delta\) by means of \(\rho\) (no. 1, Example 2); this therefore means that \(h(A_{\lambda}) \subset A_{\rho(\lambda)}'\) for all \(\lambda \in \Delta\) .

An A-homomorphism \(u: \mathbf{M} \to \mathbf{M}'\) is called graded if there exists \(\delta \in \Delta\) such that \(u\) is graded of degree \(\delta\) . If \(u \neq 0\) and every element of \(\Delta\) is cancellable, the degree \(\delta\) of \(u\) is then determined uniquely.

If \(h: \mathbf{A} \to \mathbf{A}'\) , \(h': \mathbf{A}' \to \mathbf{A}''\) are two graded homomorphisms of graded rings of type \(\Delta\) , so is \(h' \circ h: \mathbf{A} \to \mathbf{A}''\) ; for a mapping \(h: \mathbf{A} \to \mathbf{A}'\) to be a graded ring isomorphism, it is necessary and sufficient that \(h\) be bijective and that \(h\) and the inverse mapping \(h'\) be graded homomorphisms; it also suffices for this that \(h\) be a bijective graded homomorphism. Thus it is seen that graded homomorphisms can be taken as the morphisms of the species of graded ring structure of type \(\Delta\) (Set Theory, IV, § 2, no. 1).

Similarly, if \(u: \mathbf{M} \to \mathbf{M}'\) and \(u': \mathbf{M}' \to \mathbf{M}''\) are two graded homomorphisms of graded A-modules of type \(\Delta\) , of respective degrees \(\delta\) and \(\delta', u' \circ u: \mathbf{M} \to \mathbf{M}''\) is a graded homomorphism of degree \(\delta + \delta'\) . If \(\delta\) admits an inverse \(-\delta\) in \(\Delta\) and \(u: \mathbf{M} \to \mathbf{M}'\) is a bijective graded homomorphism of degree \(\delta\) , the inverse mapping \(u': \mathbf{M}' \to \mathbf{M}\) is a bijective graded homomorphism of degree \(-\delta\) . It follows as above that the graded homomorphisms of degree 0 can be taken as the morphisms of the species of graded A-module of type \(\Delta\) . But a bijective graded homomorphism \(u: \mathbf{M} \to \mathbf{N}\) of degree \(\neq 0\) is not a graded A-module isomorphism if \(\mathbf{M}\) and \(\mathbf{N}\) are non-zero and the elements of \(\Delta\) are cancellable.

Examples. (6) If M is a graded A-module and \(\mathbf{M}(\lambda_{0})\) is a graded A-module obtained by shifting (no. 2, Example 3), the Z-linear mapping of \(\mathbf{M}(\lambda_{0})\) into M which coincides with the canonical injection on each \(M_{\lambda+\lambda_{0}}\) is a graded homomorphism of degree \(\lambda_{0}\) (which is bijective when \(\Delta\) is a group).

\begin{enumerate}
\def\labelenumi{(\arabic{enumi})}
\setcounter{enumi}{6}
\tightlist
\item
  If \(a\) is a homogeneous element of degree \(\delta\) belonging to the centre of A, the homothety \(x \mapsto ax\) of any graded A-module M is a graded homomorphism of degree \(\delta\) .
\end{enumerate}

Remark (3) A graded A-module M is called a graded free A-module if there exists a basis \((m_{\iota})_{\iota \in I}\) of M consisting of homogeneous elements. Suppose it is and \(\Delta\) is a commutative group; let \(\lambda_{i}\) be the degree of \(m_{i}\) and consider for each \(\iota\) the shifted A-module \(\mathrm{A}(-\lambda_{i})\) (no. 2, Example 3); if \(e_{i}\) denotes the element l of A considered as an element of degree \(\lambda_{i}\) in \(\mathrm{A}(-\lambda_{i})\) , the A-linear mapping \(u: \bigoplus_{\iota \in I} \mathrm{A}(-\lambda_{i}) \to \mathrm{M}\) such that \(u(e_{\iota}) = m_{\iota}\) for all \(\iota\) , is a graded A-module isomorphism.

Assuming always that \(\Delta\) is a commutative group, now let N be a graded A-module, \((n_{\iota})_{\iota \in I}\) a system of homogeneous generators of N and suppose that \(n_{\iota}\) is of degree \(\mu_{\iota}\) . Then the A-linear mapping \(v: \bigoplus_{\iota \in I} A(-\mu_{\iota}) \to N\) such that \(u(e_{\iota}) = n_{\iota}\) for all \(\iota\) is a surjective graded A-module homomorphism of degree 0. If N is a finitely generated graded A-module, there is always a finite system of homogeneous generators of N and hence there is a surjective homomorphism of the above type with I finite.

